\section{Bounded Booleanization, folds, and least search}
\label{sec:booleanization}

The finite layer of the framework depends on a simple but indispensable fact:
every fixed bounded predicate can be represented by a primitive-recursive
Boolean checker whose correctness is provable in iPRA.  This section states
and proves the required form carefully.  ``Fixed'' is essential: the formula
syntax is traversed by the host compiler, not interpreted from a numerical
code inside arithmetic.

\subsection{Boolean conventions}

Use \(0\) for false and \(1\) for true.  Define primitive-recursive operations
\(\mathsf{not},\mathsf{and},\mathsf{or}\) by finite case distinction, with
outputs normalized to \(0\) or \(1\).  Equality of natural numbers has a
primitive-recursive characteristic function \(\mathsf{eq}(x,y)\).  iPRA proves
its defining facts by quantifier-free induction:
\begin{align*}
 \mathsf{eq}(x,y)=1 &\leftrightarrow x=y,\\
 \mathsf{eq}(x,y)=0 &\leftrightarrow x\neq y,\\
 \mathsf{and}(a,b)=1 &\leftrightarrow a=1\wedge b=1,\\
 \mathsf{or}(a,b)=1 &\leftrightarrow a=1\vee b=1,\\
 \mathsf{not}(a)=1 &\leftrightarrow a=0,
\end{align*}
under the Boolean normalization hypotheses where needed.

There are several possible implementations: truncated subtraction in both
directions, simultaneous predecessor recursion, or the comparison functions
already developed in the equational part of \texttt{ConciseIPRA}.  The
compiler theorem uses only the proved interface, not a particular choice.

\subsection{Quantifier-free formulas}

\begin{theorem}[Booleanization of a fixed quantifier-free formula]
\label{thm:qf-booleanization}
For every fixed quantifier-free formula \(A(\vec x)\) in the iPRA language,
one can effectively construct a primitive-recursive symbol \(b_A\) and
embedded derivations of
\[
 A(\vec x)\longrightarrow b_A(\vec x)=1,
 \qquad
 b_A(\vec x)=1\longrightarrow A(\vec x).
\]
The construction is structural in \(A\).
\end{theorem}

\begin{proof}
For an equality \(s=t\), set
\(b_A=\mathsf{eq}(s,t)\) and use the equality-checker theorem.  For falsity use
the constant-zero checker.  Conjunction and disjunction use the corresponding
Boolean operations and the induction hypotheses.

For implication \(A\to B\), both \(A\) and \(B\) are decidable because their
Boolean reflections have already been constructed.  Define
\[
 b_{A\to B}=\mathsf{or}(\mathsf{not}(b_A),b_B).
\]
If \(A\to B\) holds, decide \(A\).  In the negative branch
\(\mathsf{not}(b_A)=1\); in the positive branch obtain \(B\) and hence
\(b_B=1\).  Conversely, if the Boolean disjunction is one and \(A\) is
assumed, the negative branch contradicts the reflected truth of \(A\), while
the positive branch yields \(B\).  Every case split is on a Boolean equation
and is therefore intuitionistically valid.  No law of excluded middle for
arbitrary formulas is used.
\end{proof}

The theorem is stronger operationally than mere decidability.  It returns the
fixed checker syntax and two object derivations, which later nodes can compose.
The accompanying Lean blueprint calls this package a
\texttt{BoolReflection}.

\subsection{Bounded universal and existential folds}

Let \(b(i,\vec x)\in\{0,1\}\) be a PR Boolean predicate.  Define
\begin{align*}
 \mathsf{All}_b(0,\vec x)&=1,\\
 \mathsf{All}_b(n+1,\vec x)&=
   \mathsf{and}(\mathsf{All}_b(n,\vec x),b(n,\vec x)),\\[0.3em]
 \mathsf{Any}_b(0,\vec x)&=0,\\
 \mathsf{Any}_b(n+1,\vec x)&=
   \mathsf{or}(\mathsf{Any}_b(n,\vec x),b(n,\vec x)).
\end{align*}
These are primitive-recursive symbols by definition.

\begin{theorem}[Bounded fold reflection]
\label{thm:bounded-folds}
iPRA proves, uniformly in \(n\),
\begin{align*}
 \mathsf{All}_b(n,\vec x)=1
 &\leftrightarrow \forall i<n\; b(i,\vec x)=1,\\
 \mathsf{Any}_b(n,\vec x)=1
 &\leftrightarrow \exists i<n\; b(i,\vec x)=1,
\end{align*}
where the bounded quantifiers on the right are expanded into their standard
first-order coding or used through a bounded-formula translation.
\end{theorem}

\begin{proof}
Use quantifier-free induction on \(n\).  At zero, the universal range is empty
and the existential range is empty.  At a successor, split the bounded range
into indices below \(n\) and the final index \(n\).  The defining recursion
equations reduce the Boolean statement to a conjunction or disjunction, and
the induction hypothesis supplies the prefix statement.  All formulas used in
the induction invariant are Boolean equations and finite propositional
combinations, hence quantifier-free after the bounded side is represented by
its fold checker.
\end{proof}

The proof should be implemented with the Boolean equation as the induction
invariant, not by attempting induction on a formula containing an unbounded
quantifier.  Once reflection is proved, a bounded quantifier in a theorem can
be replaced by the equation of the corresponding fold.

\subsection{Strictly bounded formulas}

A strictly bounded formula is built from equalities and Boolean connectives,
with quantifiers of the form \(\forall i<t(\vec x)\) or
\(\exists i<t(\vec x)\).  Combining Theorems~\ref{thm:qf-booleanization} and
\ref{thm:bounded-folds} gives the following compiler.

\begin{theorem}[Booleanization of fixed strictly bounded formulas]
\label{thm:bounded-booleanization}
For every fixed strictly bounded formula \(A(\vec x)\), there is a concrete PR
symbol \(b_A\) and embedded iPRA derivations of
\[
   A(\vec x)\leftrightarrow b_A(\vec x)=1.
\]
The compiler is primitive recursive as a host transformation on the finite
formula syntax, but no universal evaluator for encoded formulas is introduced
inside iPRA.
\end{theorem}

\begin{proof}
Proceed structurally.  The quantifier-free cases use
Theorem~\ref{thm:qf-booleanization}.  For
\(\forall i<t\,B(i,\vec x)\), recursively obtain \(b_B\) and use
\(\mathsf{All}_{b_B}(t,\vec x)\); for the existential case use
\(\mathsf{Any}_{b_B}\).  Compose the reflection derivations with
Theorem~\ref{thm:bounded-folds}.  Since \(A\) is fixed, every recursively
constructed symbol has a fixed arity and is a genuine element of the
repository's inductive type \texttt{PR}.
\end{proof}

\subsection{Least bounded search}

Define a search state \((f,w)\), where \(f\in\{0,1\}\) records whether a
witness has been found and \(w\) stores the first witness.  Starting from
\((0,0)\), at index \(i\) leave the state unchanged if \(f=1\); if \(f=0\)
and \(b(i)=1\), move to \((1,i)\); otherwise remain \((0,0)\).  Pairing makes
this a primitive recursion.

Equivalently, define a sentinel-valued function \(\mu_b(n)\leq n\) by
\[
 \mu_b(n)=
 \begin{cases}
   \min\{i<n:b(i)=1\},&\text{if a witness exists},\\
   n,&\text{otherwise}.
 \end{cases}
\]

\begin{theorem}[Least bounded witness]
\label{thm:least-search}
iPRA proves all of the following:
\begin{enumerate}
\item \(\mu_b(n)\leq n\);
\item if \(\mu_b(n)<n\), then \(b(\mu_b(n))=1\);
\item if \(j<\mu_b(n)\), then \(b(j)=0\);
\item \(\mu_b(n)=n\) iff \(\mathsf{Any}_b(n)=0\).
\end{enumerate}
Thus an existence proof for a bounded decidable predicate yields an explicit
least witness without choice.
\end{theorem}

\begin{proof}
Use a single quantifier-free induction invariant recording the range bound,
the found-witness equation when the flag is one, and the prefix no-witness
fold when the flag is zero.  At the successor step, case split on the found
flag and on \(b(n)\).  The minimality clause is propagated because once the
flag becomes one the stored witness never changes.  Reflection of
\(\mathsf{Any}\) yields the final equivalence.
\end{proof}

This theorem is the constructive workhorse behind finite choice.  Whenever a
cloud is nonempty under a proved finite bound, its least member is PR.  The
important limitation is that the bound must already be present.

\subsection{Finite tables and bounded verification}

A finite table of length \(n\) can be represented by a nested pairing code or
by a PR read function generated during a bounded fold.  The framework needs
four internal theorems:
\begin{align*}
 \mathsf{read}(\mathsf{write}(T,i,x),i)&=x,\\
 j\neq i\to
 \mathsf{read}(\mathsf{write}(T,i,x),j)&=\mathsf{read}(T,j),\\
 i<n\to \mathsf{read}(T,i)&\text{ is within its declared code bound},\\
 \mathsf{AllValid}(T,n)=1&\leftrightarrow
   \forall i<n\;\Valid(\mathsf{read}(T,i),i).
\end{align*}
The last theorem is another bounded fold.  Together these facts permit a
trace invariant to speak about an entire prefix using one Boolean equation,
rather than a quantifier that would violate quantifier-free induction.

\subsection{Certificate reflection without proof reflection}

A bounded certificate checker is merely a Boolean PR predicate.  Its
reflection theorem is specific to the certificate format:
\[
 \mathsf{check}(x,c)=1\longrightarrow A(x,c).
\]
Proving this implication inside iPRA is ordinary mathematics, usually by
bounded folds and local algebra.  It must not be confused with a general
principle saying that if a numerical code is a proof, then its conclusion is
true.  PCOGs require only the former.  The checker validates a finite witness
whose semantics have been proved directly.

This distinction keeps the framework predicatively modest and makes each
certificate format independently auditable.
